\documentclass[12pt]{article}

\usepackage{ucs}
\usepackage[utf8x]{inputenc}
\usepackage[russian]{babel}
\usepackage{array,graphicx,dcolumn,multirow,abstract,hanging,fancyhdr,float}

\usepackage{amssymb}
\usepackage{amsmath}
\usepackage{indentfirst}

\newtheorem{theorem}{Теорема}

\DeclareMathOperator{\cond}{cond}
\DeclareMathOperator{\grad}{grad}
\newcommand{\norm}[1]{\left\| #1 \right\|}

\begin{document}
	\begin{flushright}
		%%.{version}.%%
	\end{flushright}
	\begin{flushright}
		%%.{date}.%%
	\end{flushright}
    \section*{Вычислительный практикум}
    \section*{Практическое занятие \#9.3. Полная проблема собственных значений. QR-алгоритм}

	\subsection*{Итерационный процесс}
	
	На 1-й итерации с помощью метода отражений или метода вращений вычисляют $\bold{QR}$-разложение матрицы
	\begin{equation}
		\bold{A}^{\left(0\right)}=\bold{Q}_1\bold{R}_1.\label{fm32}
	\end{equation}
	Затем строят матрицу $\bold{A}^{\left(1\right)}=\bold{R}_1\bold{Q}_1$. Заметим, что из равенства (\ref{fm32}) следует, что $\bold{R}_1=\bold{Q}^{-1}_1\bold{A}^{\left(0\right)}$ и поэтому $\bold{A}^{\left(1\right)}=\bold{Q}_1^{-1}\bold{A}^{\left(0\right)}\bold{Q}_1$. Таким образом, матрицы $\bold{A}^{\left(1\right)}$ и $\bold{A}^{\left(0\right)}$ подобны и поэтому имеют общий набор собственных чисел $\lambda_1,\lambda_2,\ldots,\lambda_n$. На $\left(k+1\right)$-й итерации вычисляют разложение $\bold{A}^{\left(k\right)}=\bold{Q}_{k+1}\bold{R}_{k+1}$ и строят матрицу $\bold{A}^{\left(k+1\right)}=\bold{R}_{k+1}\bold{Q}_{k+1}$.
	
	Обратим внимание, что обычной поэлементной сходимости или сходимости по норме $\bold{QR}$-алгоритм не гарантирует. Сходимость этого метода~--- это сходимость по форме последовательности приближений $\bold{A}^{\left(k\right)}$ к некоторой верхней треугольной или верхней блочно-треугольной матрице $\widetilde{\bold{A}}$, имеющей те же собственные числа, что и матрица $\bold{A}$. Сходимость по форме характеризуется сходимостью к нулю поддиагональных элементов или элементов поддиагональных блоков, в то время как наддиагональные элементы могу существенно меняться.
	
	Практически всегда $\bold{QR}$-алгоритм сходится. В общем случае, когда собственные числа кратные или комплексные, теория сходимости достаточно сложна. Ограничимся формулировкой одного из наиболее известных условий сходимости.
	
	\subsubsection*{Критерий Уилкинсона}
	
	Критерий Уилкинсона справедлив, если выполнено два условия:
	\begin{enumerate}
		\item матрица $\bold{A}$ имеет простую структуру, причем модули всех собственных значений различны
		\begin{equation}
			\left|\lambda_1\right|>\left|\lambda_2\right|>\ldots>\left|\lambda_n\right|;
		\end{equation}
		\item приведение матрицы $\bold{A}$ к диагональному виду
		\begin{equation}
			\bold{P}^{-1}\bold{AP}=\bold{D}=
			\begin{pmatrix}
				\lambda_1 & 0 & \ldots & 0 \\
				0         & \lambda_2 & \ldots & 0 \\
				\vdots & \vdots & \ddots & \vdots \\
				0 & 0 & \ldots & \lambda_n
			\end{pmatrix}
		\end{equation}
		осуществляется с помощью матрицы подобия $\bold{P}$, у которой все ведущие главные миноры отличны от нуля.
	\end{enumerate}
	
	При выполнении этих двух условий, последовательность матриц сходится по форме к верхней треугольной матрице $\widetilde{\bold{A}}$ вида
	\begin{equation}
		\widetilde{\bold{A}}=\begin{pmatrix}
			\lambda_1 & \times & \ldots & \times \\
			0         & \lambda_2 & \ldots & \times \\
			\vdots & \vdots & \ddots & \vdots \\
			0 & 0 & \ldots & \lambda_n
		\end{pmatrix},
	\end{equation}
	где знаком $\times$ обозначены элементы, в общем случае отличные от нуля.
	
	Известно, что в рассматриваемом случае элементы $a_{ij}^{\left(k\right)}$ матрицы $\bold{A}^{\left(k\right)}$, стоящие ниже главной диагонали, сходятся к нулю со скоростью геометрической прогрессии, причем
	\begin{equation}
		\left|a_{ij}^{\left(k\right)}\right|\leqslant C\left|\frac{\lambda_i}{\lambda_j}\right|^k,\quad i>j,
	\end{equation}
	то есть скорость сходимости $a_{ij}^{\left(k\right)}$ к нулю определяется значением отношения $\lambda_i$ к $\lambda_j$ (заметим, что $\dfrac{\lambda_i}{\lambda_j}<1$ при $i>j$).
	
	Если условие 2 не выполнено, то сходимость по-прежнему имеет место, но собственные значения в матрице $\widetilde{\bold{A}}$ уже не будут расположены в порядке убывания модулей.
	
	В общем случае предельная матрица $\widetilde{\bold{A}}$ --- блочно-треугольная (или получающаяся из блочно-треугольной симметричной перестановкой строк и столбцов). Наличие комплексно-сопряженных пар $\lambda_k, \overline{\lambda}_k$ собственных чисел у вещественной матрицы $\bold{A}$ не является препятствием для применения $\bold{QR}$-алгоритма. Каждой такой паре в предельной матрице будет отвечать некоторый диагональный блок --- кадратная подматрица порядка 2, собственные числа которой совпадают с $\lambda_k, \overline{\lambda}_k$.
	
	\subsection*{Ускорение метода}
	
	\subsubsection*{Сокращение арифметических операций}
	
	Каждая итерация $\bold{QR}$-алгоритма требует выполнить порядка $n^3$ арифметических операций.
	
	Чтобы уменьшить число арифметических операций, прежде чем применять $\bold{QR}$-алгоритм, следует преобразовать исходную матрицу $\bold{A}$ к форме Хессенберга
	\begin{equation}
		\bold{H}=\bold{P}^{-1}\bold{AP}=\begin{pmatrix}
			h_{11} & h_{12} & \ldots & h_{1,n-1} & h_{1n} \\
			h_{21} & h_{22} & \ldots & h_{2,n-1} & h_{2n} \\
			0      & h_{32} & \ldots & h_{3,n-1} & h_{3n} \\
			\vdots & \vdots & \ddots  & \vdots    & \vdots \\
			0      & 0      & \ldots & h_{n,n-1} & h_{nn}
		\end{pmatrix}.\label{fm35}
	\end{equation}
	Чтобы привести матрицу $\bold{A}$ к форме (\ref{fm35}), нужно выполнить последовательность преобразований с использованием элементарных матриц, таких как вращения Гивенса или отражения Хаусхолдера.
	
	Далее применяют $\bold{QR}$-алгоритм. Эффективность такого подхода обусловлена наличием следующих свойств:
	\begin{itemize}
		\item Матрицы $\bold{H}^{\left(k\right)}$, порождаемые $\bold{QR}$-алгоритмом из матрицы $\bold{H}^{\left(0\right)}$, сами являются матрицами Хессенберга, то есть для них $h_{ij}^{\left(k\right)}=0$ при $i>j+1$.
		\item Для выполнения одной итерации $\bold{QR}$-алгоритма для матрицы Хессенберга требуется число арифметических операций, пропорциональное $n^2$.
	\end{itemize}
	
	\subsubsection*{$\bold{QR}$-алгоритм со сдвигами}
	
	Также, если $\lambda_{i}\approx\lambda_{i+1}$ метод сходится очень медленно, так как например элемент $h_{i,i-1}^{\left(k\right)}$ сходится к нулю со скоростью геометрической прогрессии, знаменатель которой $q=\left|\dfrac{\lambda_i}{\lambda_{i-1}}\right|\approx 1$.
	
	В таком случае используют $\bold{QR}$-алгоритм со сдвигами. Допустим для $\lambda_i$ известно хорошее приближение $\lambda_i^{*}$, тогда собственные числа матрицы $\widetilde{\bold{H}}^{\left(k\right)}=\bold{H}^{\left(k\right)}-\lambda_i^{*}\bold{E}$ являются $\widetilde{\lambda}_j=\lambda_j-\lambda_j^{*}$, $j=1,2,\ldots,n$. В этом случае вместо отношения $\dfrac{\lambda_j}{\lambda_{i-1}}\approx 1$ скорость убывания поддиагонального элемента $\widetilde{h}^{\left(k\right)}_{i,i-1}$ определяет величина $\dfrac{\widetilde{\lambda}_i}{\widetilde{\lambda}_{i-1}}=\dfrac{\lambda_i - \lambda_{i}^{*}}{\lambda_{i-1}-\lambda_i^{*}}\approx 0$. После нескольких итераций $\bold{QR}$-алгоритма, которые практически сделают элемент $\widetilde{h}^{\left(k\right)}_{i,i-1}$ равным нулю, следует выполнить обратный сдвиг, положим $\bold{H}^{\left(k\right)}=\widetilde{\bold{H}}^{\left(k\right)}+\lambda_i^{*}\bold{E}$. После выполнения этой операции матрицы $\bold{A}$ и $\bold{H}^{\left(k\right)}$ снова имеют общий набор собственных чисел.
	
	Собственные числа найдены. Далее допустим требуется найти собственный вектор $\bold{e}_j$ матрицы $\bold{A}$, отвечающий собственному числу $\lambda_i$. Тогда сначала можно найти соответствующий собственный вектор $\bold{v}_j$ матрицы $\bold{H}$ (например, методом обратных итераций), а затем вычислить $\bold{e}_j$ по формуле $\bold{e}_j=\bold{P}\cdot\bold{v}_j$, где $\bold{P}$ --- матрицы подобия из (\ref{fm35}).
	
	\subsection*{Задания}
	
	\begin{enumerate}
		\item Реализовать $\bold{QR}$-алгоритм. При этом можно использовать метод отражений или метод вращений для вычисления $\bold{QR}$-разложения матрицы на каждой итерации.
		\item Реализовать ускорение метода.
		\item Продемонстрировать работу методов на симметричных положительно определенных матрицах.
	\end{enumerate}

    \renewcommand{\bibname}{{Список литературы}}
    \bibliographystyle{gost2008}
    \begin{thebibliography}{3}
        \bibitem{lebedeva}
        Лебедева А. В., Пакулина А. Н. Практикум по методам вычислений. Часть 1. Учебно-методическое пособие. СПб.: Санкт-Петербургский государственный университет, 2021.~156 с.
        
        \bibitem{Mis}
        Мысовских И. П. Лекции по методам вычислений: Учебное пособие. СПб.: Издательство Санкт-Петербургского университета, 1998. 472~с.
        
        \bibitem{Fadeev}
        Фаддеев Д. К., Фаддеева В. Н. Вычислительные методы линейной алгебры. М.:  Государственное издательство физико-математической литературы, 1960. 655 с.
    \end{thebibliography}
\end{document}
